\begin{tabularx}{\textwidth}{X X}
\toprule
\textbf{Előnyök} & \textbf{Hátrányok} \\
\midrule
A folyamatnak nem kell folytonos fizikai memóriaterületet kapnia. & Laptáblákra, TLB-re és hardveres címfordításra van szükség. \\
Csökkenti vagy megszünteti a külső töredezettséget. & Belső töredezettség keletkezhet az utolsó, csak részben használt lapban. \\
Nagyobb virtuális címtartományt adhat, mint a fizikai RAM mérete. & Laphiba esetén a háttértár elérése nagyon lassú a RAM-hoz képest. \\
Támogatja az igény szerinti betöltést, a megosztott könyvtárakat és a védelmet. & Rossz lapcsere vagy túl magas multiprogramozási fok thrashinget okozhat. \\
Egyszerűbbé teszi a folyamatok elkülönítését és a jogosultságok kezelését. & A címleképzés minden memóriahivatkozásnál többletmunkát jelent, amelyet gyorsítótárakkal kell csökkenteni. \\
\bottomrule
\end{tabularx}
